\documentclass[11pt]{article}
\usepackage[T1]{fontenc}
\usepackage{textcomp}
\usepackage[margin=0.75in]{geometry}
\usepackage{amsmath,amssymb,amsthm}
\usepackage{array,booktabs}
\usepackage{hyperref}
\setlength{\parindent}{0pt}
\newtheorem{theorem}{Theorem}
\theoremstyle{definition}
\newtheorem{definition}{Definition}
\title{Flow Proof Artifact: Eq-trans-subst-add}
\author{Generated by \texttt{flow doc proof}}
\date{Flow Proof Book}
\begin{document}
\maketitle
Transitive equality preserves right addition.\\[0.5em]
\textbf{Source.} leibniz — https://en.wikipedia.org/wiki/Substitution\_(logic)\\[0.5em]
\bigskip
\noindent{\large \textbf{Derived fact 1.} \textit{If a = b and b = c then a + d = c + d} \hfill equality \cdot equality \cdot transitive substitution on the right}\par
\smallskip\noindent\textit{Coordinate: equality \cdot equality \cdot transitive substitution on the right}\par
\smallskip\noindent\textit{Source: leibniz}\par
\smallskip\noindent\textit{Built on: equal terms add the same on the right, for equality on equality, equality reverses, for equality on equality}\par
\medskip
\noindent\textbf{Goal.} If a = b and b = c then a + d = c + d\par
\noindent$\displaystyle \forall a \in \mathbb{N} \forall b \in \mathbb{N} \forall c \in \mathbb{N} \forall d \in \mathbb{N}\quad a + d = c + d$\par
\medskip
\noindent
\begin{tabular}{@{} >{\raggedright\arraybackslash}p{0.06\textwidth} >{\raggedright\arraybackslash}p{0.47\textwidth} >{\raggedleft\arraybackslash}p{0.41\textwidth} @{}}
\textbf{\#} & \textbf{Proof} & \textbf{Mathematics} \\
\midrule
\textbf{1.} & We prove that transitive substitution on the right for equality on equality. & \\[0.45em]
\textbf{2.} & We split into exhaustive cases — the claim must hold in each one. & \\[0.45em]
\textbf{3.} & Case 1 (see step 2): suppose a  equals  b. & \\[0.45em]
\textbf{4.} & Case 2 (see step 2): suppose b  equals  c. & \\[0.45em]
\textbf{5.} & We invoke the derived fact governing equality on equality: equal terms add the same on the right, for equality on equality (instantiated for a, b, d). & \\[0.45em]
\textbf{6.} & We invoke the derived fact governing equality on equality: equality reverses, for equality on equality (instantiated for b, c). & \\[0.45em]
\textbf{7.} & We invoke the derived fact governing equality on equality: equal terms add the same on the right, for equality on equality (instantiated for b, c, d). & \\[0.45em]
\textbf{8.} & From step 4, step 5, step 6, and step 7, this implies a plus d equals c plus d. Together with the other cases (step 3 and step 4), the goal is discharged. Hence proven. & $\displaystyle a + d = c + d$ \\[0.45em]
\bottomrule
\end{tabular}
\label{thm:Eq-equality-transitive_substitution_on_the_right}
\par\medskip\hrule
\end{document}
